% update_reliability：生成于阅读卡+证据；引用检查通过


记忆系统的长期可靠性首先取决于更新与遗忘机制能否在不损害有用信息的前提下淘汰过时或冗余内容。现有工作从选择性遗忘、治理原语、存储/使用分离、再巩固、版本规范化与偏好更新等不同路径探索这一问题。

图结构记忆中的选择性遗忘是代表性做法。W0034 在 LongMemEval 基准上以提取式流水线评测图结构记忆（27021 个节点，候选生成预算为 5 个检索根），其选择性遗忘模块按 recency、访问频率、度中心性和 age 加权得分剪枝低分节点。作者报告，一次应用遗忘移除 9.8% 节点与 9.5% 存储字节，token F1 变化 +0.001（95% CI [-0.015, +0.016]），judged correctness 下降 1.6 点；同时图结构记忆的 token F1 为 0.417，低于扁平向量基线的 0.468（配对 bootstrap Δ=-0.050）\cite{W0034}。这说明在该设定下，图结构与剪枝策略并未带来召回优势，遗忘更多是成本控制手段而非性能提升手段。

与基于启发式得分的剪枝不同，W0199 提出 Memory Worth 作为遗忘治理的在线原语：系统在写入时记录记忆是否属于当前集合，在读取时通过任务成败反馈持续估计条件成功概率。作者称在合成环境（真实效用已知、20 个独立种子、10000 回合）中，Memory Worth 与真实效用的斯皮尔曼秩相关达到 ρ=0.89±0.02，而从不更新的系统为 ρ=0.00；在真实文本微实验（all-MiniLM-L6-v2 嵌入检索、3000 回合）中，过时记忆 MW=0.17 低于低价值阈值，专家记忆 MW=0.77。作者特别说明该度量捕获的是关联而非因果，应视为最小操作原语\cite{W0199}。

W0178 则把遗忘与信任、防投毒统一在同一预算策展框架下。CURATOR 以“净值/字节”单一评分管理经验记忆的保留、共享与信任，净值含价值项与危害项（含来源项）。作者称在 Qwen2.5-3B、Qwen2.5-7B、Phi-3.5-mini 上，仅按价值筛选或全量保留均使注入攻击成功率保持很高，而加入来源项的净值得分可将 ASR 降至零\cite{W0178}。这一结果表明，遗忘决策不能只看信息价值，还需显式建模来源可信度。

关于“该遗忘什么”，W0180 主张将存储与使用分离：归档保留源观察，查询条件视图决定哪些证据被用于回答，遗忘因而是查询局部的——事实可以继续存储但在特定答案中被抑制使用。作者在五个任务套件、四个语言模型骨干上评测 RD-Forget，报告在事实整合任务上相对更强基线的提升为 11 到 26 个百分点不等，且无遗忘配置与查询无条件配置在消融中得分最低\cite{W0180}。

W0012 提出检索驱动的再巩固（REALM）：将插入、合并与遗忘统一为单一整合阶段，把遗忘视为冗余处理。作者称 REALM 在 LoCoMo 上平均准确率 75.97%、在 LongMemEval 上 65.11%，分别比最强基线高 7.17 和 1.31 个百分点\cite{W0012}。这一思路与分离式流水线形成鲜明对比，但证据未详述冲突与信息变化的具体处理。

版本规范化处理的是跨模型记忆漂移问题。W0077 以“解释熵”刻画记忆对象被不同模型赋予不同解读的程度，主张将消歧从检索时前置到写入时。作者在 n=176 合成记忆对象、三个下游模型家族的配对准随机测试中发现，全规范化记忆相对原始记忆在情感一致性上有未校正的显著改善（+0.050，95% bootstrap CI [0.013, 0.086]，配对 t 检验 p=0.0），但多重比较校正后不再显著；事实漂移在任何校正水平下均不显著。作者明确将这些结果定性为探索性而非确认性\cite{W0077}。本文综合认为，版本规范化在概念上切中跨模型部署的痛点，但其经验证据强度尚不足以支撑强结论。

偏好更新方面，W0011 批评 MemGPT、MemoryBank、A-MEM、SCM、MemInsight 等系统默认用户行为静态、缺乏偏好漂移自适应更新，并提出融合短期滑动窗口与长期指数移动平均的偏好表征。作者称 PAMU 在 λ ∈ [0.4, 0.7] 与 β ∈ [0.8, 0.95] 范围内性能稳定\cite{W0011}。这提示偏好记忆需要同时追踪波动与趋势，但该结论的量化评测细节在所给证据中未展开【待核验】。

\subsection{非恶意可靠性}

即使没有恶意攻击，记忆系统在抽取、更新、检索与压缩各环节都可能因信息失真而不可靠。

错误传播是流水线式智能体系统的突出风险。W0688 在 4 智能体 LangGraph 流水线（Researcher→Analyst→Writer→Reviewer）上，以 FinanceBench 数据、346 条注入幻觉、gpt-4o-mini 智能体与 gpt-4o 评判进行评测。作者称幻觉检测率从 Stage 1 的 72.0% 单调下降到 Stage 4 的 50.9%；边界门控使幻觉生存率从 58.4% 降至 16.2%，而末端检查仅改善 2.3 个百分点，作者称该改善在统计上可忽略\cite{W0688}。这一对比说明，在错误产生的边界处设防比在输出末端检查更有效。

记忆幻觉可能在多个操作环节累积。W0261 提出操作级评测 HaluMem，区分记忆抽取、记忆更新与记忆问答三个任务。作者称现有记忆系统倾向于在抽取和更新阶段产生并累积幻觉，随后向问答阶段传播错误；并指出 LoCoMo、LongMemEval、PrefEval、PersonaMem 等基准采用端到端黑盒评估，难以定位幻觉发生的具体环节\cite{W0261}。本文综合认为，操作级诊断对可靠性的归因分析不可或缺，但该结论的经验基础限于 HaluMem-Medium 与 HaluMem-Long 数据集的评测范围。

检索环节的另一类失败是“证据不具可采纳性”。W0690 指出，现有评测只关注召回和最终答案正确性，无法识别检索到正确内容但因语境不匹配而不应被使用的失败模式。作者在 RHELM 与 MemOps 共 3767 个查询的 top-20 设置中报告，采用发布元数据参照后，Top-20 anchor recall 从 0.432 提升到 0.533，80% recall feasibility 从 0.237 提升到 0.311，exact similarity evaluations 下降 98.3%；在冻结 72 案例开发诊断中，仅文本验证器在 1% 必需锚点误拒上限内无法检测违规\cite{W0690}。这表明检索可采纳性需要超出语义相似度的验证信号。

来源追踪是抑制错误传播与幻觉的通用手段。W0010 提出来源感知记忆系统，存储对话事实的同时记录来源与置信度，并在读取时施加引用锁纪律：答案只能引用已打开的证据，无法组成引用集则弃权。作者报告在 LongMemEval 上达 95.60%（较此前最强系统 +0.73 分），在 LoCoMo 上达 93.60%（+1.10 分）；在八个骨干 LLM 的受控实验中，准确率仅差 3.4 分而单查询成本差约 30 倍，作者称能以约 20 倍低成本达到接近 SOTA 质量\cite{W0010}。该工作的证据强度相对较高，覆盖两个基准与多骨干模型。

从信息论视角看，所有紧凑记忆形式——KV 缓存、gist 向量、循环状态与智能体笔记——共享同一 Fano 下界约束。W0207 指出，查询无关压缩的有效每查询预算被查询分布熵 H(Q) 削减，快照式方法在查询分布变化时性能塌陷；作者还指出现有评测仅仔细测量单轮长上下文压缩，智能体实际执行的重复压缩几乎未被测量\cite{W0207}。这为理解记忆压缩失真提供了一个统一的理论参照，但其 COMPACT-Bench 的实证数据在给定证据中未提供【待核验】。

\subsection{恶意攻击与防御及其代价}

当记忆系统暴露于对抗性输入时，选择性提取与改写机制本身可能被利用。W0319 提出的 MemPoison 针对具有选择性提取与改写两阶段的 LLM 智能体记忆流水线，攻击者仅通过交互注入、无直接记忆写权限，但可借助本地白盒代理嵌入模型优化触发器。作者称 MemPoison 跨多智能体领域与记忆机制的攻击成功率最高可达 0.95，优于静态 RAG 或直接写库类基线\cite{W0319}。这揭示了一个深层矛盾：为可靠性设计的筛选流水线反而为绕过滤创造了新攻击面。

防御层面的综合策略体现在 W0076：检索时采用信任感知过滤与排序（含时间衰减），写入时通过复合信任分数门控。该工作以 guard agent 加 memory sanitization 构成防御，模型为 GPT-4o-mini-class 与 Gemini-2.0-Flash，场景为 eICU 与 MIMIC 临床 prompt 集；但给定证据未提供定量结果\cite{W0076}。本文综合判断该方向合理，但防御效果尚缺乏可核验的数字支撑【待核验】。

防御的良性用例代价常被忽视。W0312 在完全良性的 LoCoMo 对话流量上（与未防御基线同一内存后端、检索配置和冻结判断器）测量读取端 reranker 防御的副作用。作者报告该防御使核心准确率下降 4.4 点（95% CI [-9.0, 0.05]，bootstrap；McNemar p=0.064），在无攻击对话中 33.6% 的合法记忆被误隔离，单个对话最多出现 106 次误隔离，令牌开销为 2.7%；作者指出这一准确率变化处于分辨率边缘，机制性代价更显著\cite{W0312}。本文综合认为，防御评估若只报攻击成功率而不报良性成本，将系统性地高估防御可用性；W0312 与 W0178 的对照进一步表明，来源感知的评分机制可能在防投毒与保效用之间取得比纯读取端过滤更好的平衡，但这一推断需在统一基准上验证【待核验】。

\begin{figure}[t]
\centering
\begin{tikzpicture}[x=1.35cm,y=1.05cm,
  phase/.style={rectangle,rounded corners,draw,fill=blue!8,font=\scriptsize,align=center,minimum width=2.2cm},
  risk/.style={font=\tiny,color=red!60!black,align=center}]
\node[phase] (write) at (0,0) {写入\\Write};
\node[phase] (store) at (2,0) {存储\\Store};
\node[phase] (retrieve) at (4,0) {检索\\Retrieve};
\node[phase] (use) at (6,0) {使用\\Use};
\node[phase] (forget) at (2,-1.4) {遗忘/回滚\\Forget};
\node[phase] (share) at (4,-1.4) {共享/传播\\Share};
\draw[->,>=stealth,thick] (write) -- (store);
\draw[->,>=stealth,thick] (store) -- (retrieve);
\draw[->,>=stealth,thick] (retrieve) -- (use);
\draw[->,>=stealth,thick] (store) -- (forget);
\draw[->,>=stealth,thick] (store) -- (share);
\draw[->,>=stealth,dashed,thick] (use) to[bend left=25] (write);
\node[risk] at (0,0.75) {错误写入\\投毒注入};
\node[risk] at (2,0.75) {摘要失真\\来源丢失};
\node[risk] at (4,0.75) {过时信息\\检索错配};
\node[risk] at (6,0.75) {错误使用\\过度泛化};
\node[risk] at (2,-2.05) {误删/残留衍生};
\node[risk] at (4,-2.05) {跨智能体污染};
\node[risk] at (3.2,1.45) {使用反馈驱动再写入（红色为本文关注的关键风险）};
\end{tikzpicture}
\caption{记忆生命周期与关键风险示意（概念图，依据综述文献的生命周期框架
归纳\cite{W0039}；红色标注各阶段风险，详见第~\ref{sec:update}~节）。}
\label{fig:lifecycle}
\end{figure}
